Local strong convexity in the fibers requires one constant on a
neighborhood in the total space. Strong convexity of each fixed
fiber is weaker, even for a plurisubharmonic squared complex norm.

\Needspace{15\baselineskip}
\begin{example}\label{ex:fiberwise-convexity-without-local-margin}\index{Convex norm!fiberwise and local strong convexity}\index{Regularization!loss of uniform fiber margin}
On the trivial bundle $D\times\mathbb C^2$, where
$D=\{z\in\mathbb C:|z|<1/2\}$, put
\[
 Q(z,v)=\max\left\{
 |v_1|^2+|v_2|^2,
 \frac{|v_1|^2}{1-|z|^2}+|z|^2|v_2|^2
 \right\}.
\]
This is a continuous plurisubharmonic function, positive off
the zero section and complex homogeneous of degree two.
Its square root is a complex norm on every fiber, and for each
fixed $z$, the real function $Q(z,\cdot)$ is strongly convex. Nevertheless,
$Q$ is not locally strongly convex in the fibers near
$(z,v)=(0,(1,0))$ in the sense defined above.
\end{example}
\begin{proof}
Write $t=|z|^2$. Both branches are continuous, positive
semidefinite Hermitian quadratic forms in $v$, and the first
is positive definite. Their square roots satisfy the triangle
inequality and complex homogeneity; their maximum therefore
is a complex norm. In particular $Q$ is positive off zero and
continuous, also on the zero section.

The first branch is plurisubharmonic on the total space.
For the second, the locally uniformly convergent expansion
\[
 \frac{|v_1|^2}{1-|z|^2}+|z|^2|v_2|^2
   =\sum_{k=0}^{\infty}|z^kv_1|^2+|zv_2|^2
\]
expresses it as a limit of sums of squared moduli of holomorphic
functions. Thus it too is plurisubharmonic, and so is their
maximum.

For $0<t<1/4$, subtracting $t|v|^2$ gives
\[
 Q(z,v)-t|v|^2
 =\max\left\{(1-t)|v|^2,
       \left(\frac1{1-t}-t\right)|v_1|^2\right\},
\]
which is convex. At $z=0$, one has $Q(0,v)=|v|^2$.
Hence each fixed fiber has a positive strong convexity margin.

To test a margin on a neighborhood of $(0,(1,0))$, fix any
$c>0$ and choose a nonzero $z$ arbitrarily close to zero with
$t<c$. At $v=(1,0)$ the second branch is strictly larger than
the first, so it remains the maximum on a small fiber ball.
Along the real line $v=(1,s)$ in this ball,
\[
 Q(z,(1,s))-c|(1,s)|^2
   =\frac1{1-t}-c+(t-c)s^2.
\]
Its second derivative in $s$ is $2(t-c)<0$. Thus no positive $c$ makes
$Q-c|v|^2$ convex on the fiber balls in any such total-space
neighborhood. A fiber-linear change of smooth frame and its
inverse are uniformly bounded on a smaller base neighborhood,
so a positive quadratic margin in another frame would give
one in this frame. This proves the failure of the local margin.
\end{proof}
